%!TEX root = ../main.tex

\begin{abstract}

In recent years, more and more wireless sensor networks (WSN) and IoT devices are being used, and this has made their power supply a critical issue. Battery-based solutions have significant replacement and maintenance costs, especially when there are many sensor nodes or when they are hard to reach. Energy harvesting is the conversion of ambient energy, such as mechanical vibrations, heat, light, or RF signals, into usable electrical energy. It is proposed as a solution to this problem, since it allows the realisation of autonomous nodes that need almost no maintenance.

This thesis starts with an overview of the main transduction technologies used for harvesting from mechanical energy. These include piezoelectric, electromagnetic, electrostatic and triboelectric technologies, and their advantages and limitations are discussed. The thesis then focuses on the technology based on variable reluctance, called Variable Reluctance Energy Harvester (VREH). This technology works well for low-speed rotational applications, because both the magnet and the coil stay fixed, and the change in linked flux is caused only by the passage of a ferromagnetic toothed wheel. The state of the art of research on this technology is then reviewed, starting from the first prototypes built for railway monitoring, moving through single-pole topologies, and arriving at the multi-phase system with six units (MP-VREH) used as reference in this work. This system is designed to be integrated into large commercial vehicles, such as buses, at rotation speeds between 100 and 400 rpm (about 20-80 km/h).

Starting from Faraday's law, a complete electrical model of the VREH is then derived. This model gives the relations that link the extracted power to the geometric and electrical parameters of the transducer, such as the number of turns, the coil resistance, the flux difference between the aligned and unaligned positions, and the number of teeth of the wheel. From this model, the maximum power that can theoretically be transferred to the load is found, and the conditions for impedance matching are discussed as a function of the rotation speed.

The central part of the thesis focuses on the design of the power conditioning circuits for the six-phase system considered. Several rectification topologies are proposed and compared. They are first evaluated through circuit simulation, and then through an experimental test bench, which includes the design of a dedicated PCB and the calibration of the current sensors needed to accurately measure the power extracted from the system. Finally, the most suitable power management integrated circuit (PMIC) is chosen to regulate and store the harvested energy. This completes the conversion chain, from rotational mechanical energy to a stable power supply for the final load.

\end{abstract}
